% ==============================================================================
% Monografia - Áurea Santos de Oliveira
% Capítulo 7 - Considerações Finais
% ==============================================================================
\chapter{Conclusão}
\label{sec-conclusoes}

\aluno{Capítulo a ser concluído após a avaliação. Os pontos abaixo são um rascunho com o que já pode ser afirmado sobre o desenvolvimento.}

%%% Início de seção. %%%
\section{Considerações Finais}
\label{sec-conclusoes-consideracoes}

Este trabalho aplicou o \textit{framework} GraphRAG ao Código de Defesa do Consumidor, construindo uma prova de conceito que compreende a estruturação do texto da lei, a indexação em grafo de conhecimento e uma aplicação web para exploração do grafo e realização de consultas nas quatro modalidades da ferramenta, com registro detalhado de cada etapa.

Algumas lições do desenvolvimento já podem ser registradas:

\begin{itemize}
    \item \textbf{A preparação do corpus importa.} Estruturar a lei em um registro por artigo, com a hierarquia e a situação normativa inseridas em cada trecho, fez com que cada unidade de texto correspondesse a um artigo completo e permitiu que o grafo registrasse corretamente localização, vetos e revogações.
    \item \textbf{A qualidade do grafo depende da padronização.} Como o GraphRAG funde entidades pelo nome e pelo tipo, as definições dos tipos e as regras de nomeação inseridas no \textit{prompt} de extração foram decisivas para evitar nós duplicados.
    \item \textbf{O custo é dominado pela extração e pelas modalidades com muitas chamadas.} A indexação custou várias vezes a estimativa feita a partir do tamanho do corpus, e as buscas global e DRIFT custam dezenas de vezes uma busca local, conforme a configuração.
    \item \textbf{Usar a biblioteca exige ler o código.} Defeitos e comportamentos não documentados -- como a busca básica que devolve resposta vazia quando usada programaticamente, as buscas locais da DRIFT que não usam o grafo com a configuração padrão e o ajuste de \textit{prompts} que gera todos os exemplos a partir do mesmo trecho -- afetam diretamente a comparação entre as modalidades e só foram identificados pela inspeção do código e dos dados intermediários. A aplicação desenvolvida, ao expor cada etapa das consultas, foi o que tornou possível identificá-los.
\end{itemize}

\pendente{[Relacionar os resultados da avaliação aos objetivos da Seção~\ref{sec-intro-obj}; limitações; dificuldades; experiência adquirida.]}

%%% Início de seção. %%%
\section{Trabalhos Futuros}
\label{sec-conclusoes-trabalhosfuturos}

\pendente{[Rascunho de possibilidades:]}

\begin{itemize}
    \item ampliar o corpus com outras leis relacionadas (por exemplo, a Lei da Ação Civil Pública e regulamentações do CDC) e com jurisprudência, aproveitando o \textit{parser} genérico e as remissões entre normas já representadas no grafo;
    \item explorar a extração de alegações (\textit{claims}) e a seleção dinâmica de comunidades;
    \item comparar provedores e modelos de linguagem, inclusive modelos locais, mantendo o restante da configuração;
    \item reportar à equipe do GraphRAG os defeitos identificados, com as correções propostas;
    \item avaliar automaticamente as respostas com um LLM como juiz, comparando com a avaliação manual.
\end{itemize}
